\section{Constructing Reasoning Trees from Episode-Annotated Traces}
\label{app:reasoning_tree_construction}

\paragraph{Method and inputs.}
We adapt LCoT2Tree~\citep{jiang-etal-2025-makes} to characterize the
structure of extracted reasoning. Each trace is represented by an ordered
reasoning sketch and a tree that places passages at the corresponding
sketch steps. Our adaptation uses semantic episode boundaries, distinguishes
executing a step from referring to it, and retains brief computations and
checks as leaves.

We analyze one completed forced-reasoning trace per model and problem on
MATH: GPT-6 Astra with think-here and low native effort, GPT-5.6 Sol with
maximal deliberation, Claude Opus~4.8, and Claude Sonnet~5. Both correct and
incorrect responses are included; when several completed traces are
available, we select the first in generation order. Extracted reasoning traces from
successive tool calls are concatenated chronologically. The separately
emitted final answer and the reference answer are excluded from the judge's
input.

\paragraph{Episode-based segmentation.}
Recent models such as Astra and Sol can express reasoning more compactly,
with changes in reasoning activity that are not always marked by explicit
transition phrases. Astra, in particular, rarely uses cues such as
``Wait'' or ``Let me verify,'' making keyword-based segmentation insufficient.
We therefore identify thought boundaries using the semantic episode
annotations described in Appendix~\ref{app:llm_judge}.
A \emph{thought} consists of consecutive annotated units.
A new thought begins at each \textsc{Explore}, \textsc{Verify}, or
\textsc{Monitor} unit; other units continue the current thought.
Consecutive units in any of these three categories begin separate thoughts.
The same rule applies to all models.


\paragraph{Sketch and step assignment.}
We use DeepSeek-V4-Flash with native reasoning disabled
to extract an ordered sketch, assign thoughts to its steps, and classify
transitions. Sketch steps describe the reasoning within an individual
trace; their indices do not align mathematical content across models.

For each thought $i$, the judge distinguishes the steps it carries out,
derives, checks, or revises ($W_i$) from those whose results it merely cites
or uses ($R_i$). This distinction prevents a reference to an earlier result
from being counted as a return to that step. Assignments must correspond to
steps and thoughts present in the trace's representation. A thought may
execute several steps or none.

\paragraph{Tree construction.}
Starting from an artificial root at level~0, thoughts are incorporated in
textual order. For a thought with nonempty $W_i$, the tree follows the
current branch back to the nearest ancestor below its first assigned step,
then adds a chain of nodes at the assigned steps in their given order.
Subsequent thoughts continue from the end of this chain. Revisiting a step
creates a new node, preserving repeated work and alternative paths.
Transitions are labeled as continuation, exploration, backtracking, or
validation; these labels describe the transition but do not determine
parent placement.

Brief computations and checks may receive only reference assignments.
When $W_i$ is empty, $R_i$ is nonempty, and the thought contains an
\textsc{Implement} or \textsc{Verify} unit, we retain it as a
\emph{short-check leaf}. Among its referenced steps already represented in
the tree, we choose the highest step and its most recent occurrence. The
leaf is attached to that occurrence's parent at the same level, while the
main reasoning branch continues from its previous position. No leaf is
added if none of the referenced steps is represented. These leaves
contribute to tree width and size.

\paragraph{Metrics and interpretation.}
Width $p$ is the maximum number of non-root nodes at a sketch-step level;
depth $q$ is the largest occupied step index; and size $N$ is the total
number of non-root nodes. Depth thus measures progression through the
sketch, rather than root-to-node path length. Nodes are displayed by sketch
level; horizontal spacing has no temporal or computational interpretation.

These metrics describe the organization of visible reasoning and depend on
segmentation and step assignment. Smaller trees indicate fewer represented
nodes, but do not by themselves establish fewer internal computations or
fewer necessary proof steps.

Table~\ref{tab:reasoning-tree-structure} reports medians and quartiles;
Table~\ref{tab:reasoning-tree-means} gives means over the same traces.
Astra has smaller mean width and size, but greater mean depth; the
comparison of similar depths in the main text refers to medians.

\begin{table}[ht]
\centering
\caption{Mean reasoning-tree width, depth, and size on MATH, using the
same traces and metric definitions as Table~\ref{tab:reasoning-tree-structure}.}
\label{tab:reasoning-tree-means}
\small
\begin{tabular}{lrrr}
\toprule
Model & Width $p$ & Depth $q$ & Nodes $N$ \\
\midrule
GPT-6 Astra & 7.3 & 14.9 & 41.0 \\
GPT-5.6 Sol & 16.7 & 11.2 & 65.9 \\
Claude Opus 4.8 & 30.6 & 11.9 & 144.2 \\
Claude Sonnet 5 & 44.3 & 12.4 & 171.4 \\
\bottomrule
\end{tabular}
\end{table}

\paragraph{Step-assignment prompt.}
The prompt used to assign thoughts to reasoning-sketch steps is shown below.

\begin{quote}
\footnotesize\ttfamily\raggedright
Your task is to match each reasoning thought from List B to corresponding step number(s) in the List A, and for every match say whether the thought WORKS ON that step or only REFERS TO it. Follow the following process:\par
1. FIRST UNDERSTAND LIST B:\par
   - For each thought in List B, identify if it describes some SPECIFIC CALCULATION PROCESSes (mathematical operation, logical transformation, or data manipulation)\par
   - Ignore the description that only state conclusions, concepts without showing the actual processing detail\par
2. THEN MATCH TO LIST A:\par
   - For each thought from List B, find all steps in List A that:\par
     * Show the same underlying calculation (even with different numbers/words)\par
     * Represent the partial or same reasoning process\par
   - Ignore superficial wording differences - focus on logical equivalence\par
   - Put a step under "works" when the thought actually carries out, derives, checks or revises that step's calculation.\par
   - Put a step under "refers" when the thought only mentions, restates, or uses that step's result as given, without redoing its calculation.\par
3. OUTPUT REQUIREMENTS:\par
   - Return ALL plausible matches where computational processes align\par
   - A thought that performs no calculation at all gets "works": [] (it may still have "refers")\par
   - Multiple matches are encouraged when justified\par
   - Every thought of List B must appear as a key\par
   - Maintain strict JSON format\par
Input:\par
- List A (Detailed Steps): \par
\textless{}list\_a\textgreater{}\par
\{reasoning\_step\}\par
\textless{}/list\_a\textgreater{}\par
- List B (Reasoning Thoughts): \par
\textless{}list\_b\textgreater{}\par
\{thoughts\}\par
\textless{}/list\_b\textgreater{}\par
Output Format (strict JSON):\par
```json\par
\{\par
  "B0": \{"works": ["A1"], "refers": []\},\par
  "B1": \{"works": ["A3"], "refers": ["A1"]\},\par
  "B2": \{"works": ["A1", "A4"], "refers": []\},\par
  ...\par
\}```\par
Please match the reasoning thoughts in List B to step in the List A.\par
\end{quote}
